\cvsection{Experience}

\begin{cventries}
  \cventry
    {AI Research Engineering Intern}
    {INRIA}
    {Bordeaux, France}
    {Mar. 2026 -- Sept. 2026}
    {
      \begin{cvitems}
        \item {Independently designed and built a working application and AI agent to help researchers find scientific papers, with guidance from my supervisors.}
        \item {Developed a Python/FastAPI backend with REST APIs documented using OpenAPI and MongoDB persistence.}
        \item {Built a React/Next.js interface for interacting with the agent and its tools, including user selection of messages to include in the model context.}
        \item {Implemented a multi-step agent without an orchestration framework to understand its underlying mechanisms, integrating the OpenAI API, prompts, tool execution, state and memory management.}
        \item {Implemented a RAG pipeline with document chunking, embeddings and semantic retrieval using ChromaDB.}
        \item {Automated online scientific-paper metadata collection in Python and delivered cron task results through webhooks.}
        \item {Containerized components with Docker and set up a GitHub CI/CD pipeline. Managed a web application using LLMs installed locally on a Mac Studio, including model installation and configuration.}
        \item {Used vLLM to serve models on HPC infrastructure with NVIDIA GPUs. Compared Qwen, Gemma and GLM and quantization options using Slurm, measuring latency, TTFT, response quality and context length.}
        \item {Defined the experimental protocol, analyzed results and worked on writing a scientific article comparing language models.}
      \end{cvitems}
    }

  \cventry
    {Reinforcement Learning Research Intern}
    {Karunya University}
    {Coimbatore, India}
    {June 2025 -- Sept. 2025}
    {
      \begin{cvitems}
        \item {Developed a reinforcement learning model and its training environment for autonomous navigation of a hospital robot.}
        \item {Deployed and optimized the software and model on a Raspberry Pi for real-world testing.}
      \end{cvitems}
    }
\end{cventries}
